\section{The Capability Lattice}
\label{sec:lattice}

The foundational paper defines a capability as ``something an agent or
coalition of agents could do or be'' and the capability space $\G$ as the
universe of all such capabilities. This section makes both concrete by
constructing $\G$ as a finite weighted poset
\citep{davey2002introduction} discovered through communication and observation.

% ---------------------------------------------------------------------------
% Capability Descriptions and Benchmarks
% ---------------------------------------------------------------------------

\begin{definition}[Capability Description]
\label{def:capability_description}
A \emph{capability description} $d$ is a finite expression in a shared
    language $L$ specifying what an agent can do, under what conditions. The
    language $L$ must be expressive enough that agents can describe their
    capabilities to each other and that the GFM actor can parse descriptions
    into structured objects.
\end{definition}

Capability descriptions are the atomic currency of the lattice, analogous to
attributes in formal concept analysis \citep{ganter1999formal}. ``Can run a
mile in under 8 minutes,'' ``can write a sorting algorithm in Python,'' and
``can perform cardiac surgery with an anesthesiologist'' are all capability
descriptions. The shared language $L$ need not be natural language; any
structured representation (a formal specification language, a capability
ontology \citep{gruber1993translation}, a feature vector with semantic labels)
suffices, provided agents can
produce and consume expressions in $L$.

\begin{definition}[Benchmark]
\label{def:benchmark}
A \emph{benchmark} $B(d)$ for a capability description $d$ is a procedure
    that maps agents to a bounded numerical scale $[0, \smax]$, measuring the
    degree to which an agent possesses capability $d$. A benchmark must be:
\begin{enumerate}
\item[\textbf{(B1)}] \emph{Communicable:} the procedure can be described in $L$.
\item[\textbf{(B2)}] \emph{Repeatable:} applying $B(d)$ to the same agent
    under similar conditions yields consistent results (bounded variance).
\item[\textbf{(B3)}] \emph{Bounded:} $B(d)$ has finite range $[0, \smax]$ with
    $\smax < \infty$.
\end{enumerate}
\end{definition}

Every capability known to the GFM actor carries at least binary weight. If the
actor knows a capability exists (through communication or observation) but has
no graded benchmark, it assigns the default benchmark $\smax = 1$ (pass/fail),
giving $w(d) = \log_2(2) = 1$ bit: the actor knows one bit of information
about the capability (can/cannot). A marathon time (measured in seconds, $\smax
= 21600$), a programming test score (measured in problems solved, $\smax =
100$), and a binary certification (pass/fail, $\smax = 1$) are all valid
benchmarks at different granularities. Developing a finer benchmark ($\smax >
1$) can only increase $w(d)$ beyond the binary default, so benchmark
development is always $\volL$-non-decreasing. The zero-weight gap is confined
to capabilities the actor does not know about at all, not capabilities it knows
about but has not measured precisely.

% ---------------------------------------------------------------------------
% Subsumption and the Partial Order
% ---------------------------------------------------------------------------

\begin{definition}[Subsumption]
\label{def:subsumption}
Capability $d_1$ \emph{subsumes} $d_2$ (written $d_2 \leq d_1$) if every
    agent that achieves $d_1$ at any positive benchmark level necessarily
    achieves $d_2$ at some positive level. The relation $\leq$ is a partial
    order on the set of benchmarked capability descriptions.
\end{definition}

``Can run a marathon'' subsumes ``can run a mile'': any agent that completes a
marathon has necessarily run a mile. ``Can build a web application'' subsumes
``can write a function'': building an application requires writing functions.
The subsumption relation captures the hierarchical structure of capabilities
that a flat enumeration would miss.

Not all pairs of capabilities are comparable. ``Can run a marathon'' and ``can
play piano'' are incomparable: neither subsumes the other. The partial order
captures exactly this: capabilities that share no prerequisite structure occupy
incomparable positions in the poset.

% ---------------------------------------------------------------------------
% The Weight Function
% ---------------------------------------------------------------------------

\begin{definition}[Weight Function]
\label{def:weight}
The \emph{weight} of a benchmarked capability $d$ is:
\begin{equation}
\label{eq:weight}
w(d) = \log_2(1 + \smax(d))
\end{equation}
where $\smax(d)$ is the number of distinguishable levels in $d$'s benchmark
    scale.
\end{definition}

The weight is the information content of the benchmark in bits
\citep{cover2006elements}. A binary
capability (pass/fail, $\smax = 1$) has weight $w = 1$ bit. A capability
measured on a 1000-level scale has weight $w \approx 10$ bits. This choice
connects to the empowerment interpretation of the foundational paper: the
information content of a benchmark is analogous to the channel capacity of the
capability's action-state mapping.

The log weight ensures $w(d) > 0$ for any non-trivial benchmark ($\smax \geq
1$), satisfying the spirit of axiom M5 (non-triviality). It also ensures that
weight grows sublinearly with benchmark granularity: doubling the precision of
a benchmark does not double its contribution to $\volL(G)$, reflecting the
diminishing marginal value of finer measurement.

% ---------------------------------------------------------------------------
% The Lattice and Its Measure
% ---------------------------------------------------------------------------

\begin{definition}[Capability Lattice]
\label{def:capability_lattice}
The \emph{capability lattice} $(\C, \leq)$ is the finite partially ordered
    set of all benchmarked capability descriptions known to the GFM actor,
    ordered by subsumption (Definition~\ref{def:subsumption}), augmented with
    a bottom element $\bot$ (the null capability, $w(\bot) = 0$).
\end{definition}

\begin{definition}[Independent Weight]
\label{def:independent_weight}
The \emph{independent weight} of a capability $d \in \C$ is defined via
    M\"{o}bius inversion \citep{rota1964foundations} on $(\C, \leq)$:
\begin{equation}
\label{eq:independent_weight}
\iw(d) = \sum_{d' \leq d} \mu(d', d) \cdot w(d')
\end{equation}
where $\mu$ is the M\"{o}bius function of the poset $(\C, \leq)$, defined
    recursively by $\mu(d, d) = 1$ and
    $\mu(d', d) = -\sum_{d' \leq d'' < d} \mu(d', d'')$ for $d' < d$.
\end{definition}

The independent weight is the marginal contribution of $d$ beyond what its
subsumees already provide. The M\"{o}bius function handles the
inclusion-exclusion correctly even when the lattice has branching and shared
descendants: in a diamond $x < y, z < t$, the function accounts for $x$ being
reachable from $t$ through both $y$ and $z$, avoiding the double-counting that
a cover-only subtraction would produce.

For a chain (no branching), the M\"{o}bius function reduces to the familiar
cover subtraction: if ``marathon'' ($w = 14.4$) subsumes only ``mile'' ($w =
9.2$), then $\iw(\text{marathon}) = 14.4 - 9.2 = 5.2$ and $\iw(\text{mile})
= 9.2$. Together they contribute $9.2 + 5.2 = 14.4 = w(\text{marathon})$: the
marathon's full weight, without double-counting.

\begin{definition}[Lattice Measure $\volL$]
\label{def:lattice_measure}
For a set of capabilities $G \subseteq \C$, define the \emph{lattice measure}:
\begin{equation}
\label{eq:lattice_measure}
\volL(G) = \sum_{d \in \dc G} \iw(d)
\end{equation}
where $\dc G = \{d' \in \C : d' \leq d \text{ for some } d \in G\}$ is the
    \emph{downward closure} of $G$ under $\leq$.
\end{definition}

The downward closure ensures that including a capability in $G$ automatically
includes everything it subsumes. The lattice measure sums the independent
weights over this closure, giving each capability credit for exactly its
marginal contribution.

\begin{lemma}[Non-Negativity of Independent Weights]
\label{lem:nonneg_iw}
If $(\C, \leq)$ is a \emph{forest} (every element has at most one cover),
    then $\iw(d) \geq 0$ for all $d$ whenever $w(d) \geq w(d')$ for $d'
    \lessdot d$. For general lattices with branching, the sufficient condition
    is that $w$ is \emph{M\"{o}bius-non-negative}: $\sum_{d' \leq d}
    \mu(d', d) \cdot w(d') \geq 0$ for all $d$.
\end{lemma}

\begin{proof}
On a forest (tree-structured poset), each element $d$ has at most one cover
$d'$, so $\mu(d', d) = -1$ for the unique cover and $\mu(d, d) = 1$, giving
$\iw(d) = w(d) - w(d')$. The condition $w(d) \geq w(d')$ is sufficient.

On a general lattice, the M\"{o}bius function can produce alternating-sign
terms. In a diamond $\bot < y, z < t$, $\mu(\bot, t) = 1$ (not $0$), so
$\iw(t) = w(t) - w(y) - w(z) + w(\bot) = w(t) - w(y) - w(z)$. This can be
negative even when $w(t) \geq \max(w(y), w(z))$. The condition $\iw(d) \geq
0$ for all $d$ is therefore a constraint on the relationship between the weight
function and the lattice topology, verified at lattice construction time.
\end{proof}

\paragraph{Practical scope of the constraint.} Capability lattices in practice
are predominantly tree-structured: most capabilities have a single primary
prerequisite chain (mile $\to$ marathon, function $\to$ webapp). Diamond
structures arise from cooperative capabilities that subsume multiple
prerequisites (musicapp subsumes both webapp and perform). For these nodes, the
GFM actor verifies $\iw(d) \geq 0$ at insertion time. If the check fails (the
supersumer's benchmark is too coarse relative to the combined prerequisites),
the actor can either refine the benchmark to increase $w(d)$ or decline to add
the subsumption edges, treating the capability as lattice-disjoint from its
would-be prerequisites. The worked example in Section~\ref{sec:example}
demonstrates a lattice that passes the check for all nodes including its
diamond.

\paragraph{Cooperative capabilities.} A cooperative capability $d_{\mathrm{coop}}$
is a lattice node whose benchmark conditions require the participation of
multiple agents. ``Can perform cardiac surgery'' requires both a surgeon and an
anesthesiologist; its benchmark might measure patient outcomes on a scale from
0 to 100. Cooperative capabilities enter the lattice through the same channels
as individual capabilities (Section~\ref{sec:dual_channel}) and may subsume
individual capabilities: ``can perform cardiac surgery'' subsumes both ``can
operate'' and ``can manage anesthesia.''

% ---------------------------------------------------------------------------
% Dual-Channel Discovery
% ---------------------------------------------------------------------------

\subsection{Dual-Channel Capability Discovery}
\label{sec:dual_channel}

The lattice is populated from two independent channels, providing robustness
against agents that misrepresent or conceal their capabilities.

\paragraph{Communication channel.} Agent $a_k$ tells the GFM actor what it can
do by transmitting capability descriptions with claimed benchmark scores. The
claim enters the lattice weighted by trust:
\begin{equation}
\label{eq:w_comm}
\wcomm(d, k) = T_k \cdot w(d)
\end{equation}
where $T_k \in [0, 1]$ is the trust factor from the foundational paper's
observable goal model. This is the channel that the original $R_k$ signal
operates on.

\paragraph{Observation channel.} The GFM actor observes agent $a_k$
demonstrating capability $d$ (e.g., $a_k$ successfully performs an action that
requires $d$). The observation enters the lattice weighted by observational
confidence:
\begin{equation}
\label{eq:w_obs}
\wobs(d, k) = \Cobs(d, k) \cdot w(d)
\end{equation}
where $\Cobs(d, k) \in [0, 1]$ reflects the actor's certainty that the
observed behavior constitutes demonstration of $d$.

\paragraph{Divergence as trust signal.} When the two channels persistently
disagree for the same capability:
\begin{equation}
\label{eq:divergence}
\bigl|\wcomm(d, k) - \wobs(d, k)\bigr| > \theta_d
\end{equation}
the GFM actor registers a prediction residual in the trust model. Three
patterns are distinguishable:

\begin{enumerate}
\item \emph{Overclaiming} ($\wcomm \gg \wobs$): agent claims capabilities it
  never demonstrates. Trust $T_k$ decreases. Indicates deception or
  self-delusion.
\item \emph{Underclaiming} ($\wcomm \ll \wobs$): agent demonstrates
  capabilities it does not report. Trust $T_k$ is not penalized (the agent is
  not deceptive), but the actor expands its model of $a_k$. Indicates modesty
  or strategic concealment.
\item \emph{Consistent} ($\wcomm \approx \wobs$): claims match behavior.
  Trust $T_k$ is maintained or increases. Indicates reliability.
\end{enumerate}

\paragraph{Effective weight.} A capability's effective weight for agent $a_k$
is the maximum of the two channel estimates:
\begin{equation}
\label{eq:w_eff}
\weff(d, k) = \max\bigl(\wcomm(d, k),\; \wobs(d, k)\bigr)
\end{equation}
Taking the maximum ensures that an agent's capabilities are credited through
whichever channel provides stronger evidence. An agent that refuses to
communicate is still credited for demonstrated capabilities. An agent that
overclaims may temporarily retain the higher communicated weight, but the
divergence signal erodes $T_k$ over time, eventually driving $\wcomm(d, k)$
below $\wobs(d, k)$ and correcting the effective weight toward the observed
level. The correction is eventual, not instantaneous: at any given moment,
$\weff$ may reflect an overclaim that has not yet been fully discounted.

The dual-channel approach does not increase the computational complexity of
$\volL$: the lattice structure is the same regardless of how nodes are
populated. It increases the \emph{quality} of the lattice by providing
independent evidence for each capability claim, and it makes the trust model
operate on structured capability data rather than ternary signals.
